\ifodd\value{diceNo}

\spell{Sky-Call}% Name
  {Distant, Detailed}% Enhancements
  {Wane}% Action
  {Earth}% Spheres
  {\roll{Wits}{Caving}}% Resist with
  {The caster calls to a distant stone ceiling or wall, commanding it to crack and splint, and send a hunk of rock onto someone's head.
    If the target fails to notice the hunk of rock, they receive \showDam\ Damage}% Description
  {
    The ceiling generally remains intact, despite the small missing chunk, although repeated castings will inevitably bring down an entire floor of a castle, and \emph{might} cause a cave-in when underground.

    The spell can target ice as easily as any other type of solid material, but cannot fracture wood.
    Dense rock may give a penalty to the spell's \gls{tn}, at the \gls{gm}'s discretion.

    A failed roll generally indicates that the crack formed too slowly, and made too much noise, signaling to anyone below that rock might fall at any moment.
    Rolls with a high failure-margin might indicate that the rock failed to move entirely, and that anyone below the ceiling clearly understood that nothing could move or break it.
  }

\else

\spell{Undress the Hill}% Name
  {Detailed, Duplicated}% Enhancements
  {Wane}% Action
  {Earth}% Spheres
  {solidity of the earth}% Resist with
  {The caster compliments the ground, comparing it to a palace.
  A defined shape of \arabic{spellTargets}~\glspl{step} in total cracks, crumbles, or turns to sludge, leaving only the shape which the caster described standing.
  Anyone standing on the collapsing exterior receives a -\arabic{spellPlusTwo}~Penalty to \glspl{action} involving balance}% Description
  {
	Destroying earth on a slope will send it sliding downhill.
    Targeting soft earth during the rain has \tn[7], but the unblemished shape which remains will not have enough stability to remain upright if everything around it has washed downhill.

    \Gls{casting} upon a mountain's stone top allows the caster to quickly carve out a properly defined space which remains standing once the surrounding rocks tumble downhill, but the \gls{tn} rises to 12 or more.
  }

\fi
\stepcounter{diceNo}
